\documentclass[12pt]{article}
\usepackage{amsmath, amssymb, geometry, enumitem, booktabs}
\geometry{margin=1in}
\setlength{\parskip}{0.6em}
\setlength{\parindent}{0pt}

\begin{document}

\begin{center}
\Large \textbf{ECON 300 -- Labour Economics} \\[0.7em]
\large \textbf{Practice Problem Set 5}\\[0.3em]

\end{center}

\vspace{1em}


\begin{enumerate}[label=\textbf{Question \arabic*.}]

\item A country admits 320{,}000 immigrants in a year. Of these:
\begin{itemize}
    \item 160{,}000 are economic immigrants,
    \item 96{,}000 are family-class immigrants,
    \item 64{,}000 are refugees or other humanitarian immigrants.
\end{itemize}

Answer the following:
\begin{enumerate}[label=(\alph*)]
    \item What percentage of total immigrants are economic immigrants?
    \item What percentage are family-class immigrants?
    \item What percentage are refugees or humanitarian immigrants?
    \item If next year the number of economic immigrants rises by 20\% while the other two categories remain unchanged, what will be the new total number of immigrants?
    \item Under the new total, what share of all immigrants will be economic immigrants?
\end{enumerate}

\item In a local labour market, the annual labour demand and labour supply for low-skilled labour are given by:

\[
W = 32 - 0.02N^D
\]

\[
W = 8 + 0.01N^S
\]

where \(W\) is the hourly wage and \(N\) is the number of workers.

\begin{enumerate}[label=(\alph*)]
    \item Find the initial equilibrium wage and employment.
    \item Suppose immigration increases labour supply by 120 workers at every wage, so the new supply curve becomes:
    \[
    W = 8 + 0.01(N^S - 120)
    \]
    Find the new equilibrium wage and employment.
    \item Calculate the change in wage and the change in employment caused by immigration.
    \item Briefly interpret the result.
\end{enumerate}

\item Suppose immigrants and native-born workers are imperfect substitutes. The labour demand for native-born workers is:

\[
W = 40 - 0.01N_n + 0.005N_i
\]

where:
\begin{itemize}
    \item \(W\) = native-born wage,
    \item \(N_n\) = number of native-born workers,
    \item \(N_i\) = number of immigrant workers.
\end{itemize}

Initially, \(N_i = 200\), and native-born labour supply is:

\[
W = 10 + 0.02N_n
\]

\begin{enumerate}[label=(\alph*)]
    \item Find the equilibrium wage and employment of native-born workers when \(N_i = 200\).
    \item Suppose immigration rises so that \(N_i = 300\). Find the new equilibrium wage and employment of native-born workers.
    \item Explain whether immigration helps or hurts native-born workers in this example.
\end{enumerate}

\item A 25-year-old worker compares the present value of lifetime earnings in Country A and Country B. The worker expects to work for 35 more years.

\begin{itemize}
    \item Annual earnings in Country A: \$52{,}000
    \item Annual earnings in Country B: \$58{,}000
    \item Annual non-monetary migration cost equivalent: \$2{,}000
    \item One-time moving cost: \$15{,}000
    \item Discount rate: 4\%
\end{itemize}

Assume annual earnings are constant over time.

\begin{enumerate}[label=(\alph*)]
    \item Compute the annual net earnings advantage of Country B over Country A.
    \item Compute the present value of this annual earnings advantage over 35 years using the annuity formula:
    \[
    PV = A\left(\frac{1-(1+r)^{-T}}{r}\right)
    \]
    \item Subtract the one-time moving cost and determine whether migration is worthwhile.
    \item If the moving cost rises to \$30{,}000, does the decision change?
\end{enumerate}

\item Consider the following yearly earnings profile of an immigrant relative to a comparable native-born worker:

\begin{center}
\begin{tabular}{ccc}
\toprule
Years Since Arrival & Immigrant Earnings (\$) & Native-born Earnings (\$) \\
\midrule
0 & 32{,}000 & 44{,}000 \\
5 & 39{,}000 & 47{,}000 \\
10 & 46{,}000 & 50{,}000 \\
15 & 52{,}000 & 54{,}000 \\
\bottomrule
\end{tabular}
\end{center}

\begin{enumerate}[label=(\alph*)]
    \item Compute the earnings gap at arrival.
    \item Compute the earnings gap after 5, 10, and 15 years.
    \item By how much has the immigrant closed the earnings gap after 15 years?
    \item What percentage of the initial gap has been closed after 15 years?
\end{enumerate}

\item In a point system, applicants are scored as follows:
\begin{itemize}
    \item Education: up to 25 points
    \item Official language ability: up to 28 points
    \item Work experience: up to 15 points
    \item Age: up to 10 points
    \item Arranged employment: up to 10 points
    \item Adaptability: up to 10 points
\end{itemize}

An applicant receives:
\begin{itemize}
    \item Education: 21
    \item Language: 20
    \item Work experience: 11
    \item Age: 8
    \item Arranged employment: 0
    \item Adaptability: 7
\end{itemize}

\begin{enumerate}[label=(\alph*)]
    \item Calculate the total points.
    \item If the passing score is 67, does the applicant qualify?
    \item If the applicant improves language points by 5, what is the new total?
    \item What is the minimum number of additional points needed to reach 67 if the applicant initially fails?
\end{enumerate}

\newpage
\item In a labour market:
\begin{itemize}
    \item Total working-age population = 2{,}400{,}000
    \item Employed = 1{,}560{,}000
    \item Unemployed = 120{,}000
\end{itemize}

\begin{enumerate}[label=(\alph*)]
    \item Calculate the labour force.
    \item Calculate the unemployment rate.
    \item Calculate the labour force participation rate.
    \item Calculate the employment rate.
\end{enumerate}

\item Suppose a country has:
\begin{itemize}
    \item Labour force = 1{,}800{,}000
    \item Unemployment rate = 7.5\%
\end{itemize}

\begin{enumerate}[label=(\alph*)]
    \item Calculate the number of unemployed workers.
    \item Calculate the number of employed workers.
    \item If employment rises by 45{,}000 and unemployment falls by the same amount, what is the new unemployment rate?
\end{enumerate}

\item In a given month, the following labour market transitions occur:
\begin{itemize}
    \item 210{,}000 unemployed workers find jobs
    \item 180{,}000 employed workers lose or leave jobs and begin searching
\end{itemize}

\begin{enumerate}[label=(\alph*)]
    \item Calculate the net monthly change in unemployment.
    \item If the initial stock of unemployed workers is 960{,}000, what is the stock of unemployed workers at the end of the month?
    \item If the labour force is 16{,}000{,}000, what is the unemployment rate at the end of the month?
\end{enumerate}
\newpage
\item In a labour market, the unemployment rate can be approximated by:

\[
UR = \text{Incidence} \times \text{Duration}
\]

Suppose:
\begin{itemize}
    \item Region A has incidence = 2.5\% per month and average duration = 3 months
    \item Region B has incidence = 1.5\% per month and average duration = 5 months
\end{itemize}

\begin{enumerate}[label=(\alph*)]
    \item Calculate the unemployment rate for Region A.
    \item Calculate the unemployment rate for Region B.
    \item Which region has the higher unemployment rate?
    \item Explain in words how the same unemployment rate could arise from different combinations of incidence and duration.
\end{enumerate}

\item The following information is given for two age groups:

\begin{center}
\begin{tabular}{lccc}
\toprule
Group & Unemployment Rate & Incidence Rate & Duration (months) \\
\midrule
Youth (15--24) & 12.0\% & 6.0\% & 2.0 \\
Adults (25--54) & 5.6\% & 2.0\% & 2.8 \\
\bottomrule
\end{tabular}
\end{center}

\begin{enumerate}[label=(\alph*)]
    \item Verify approximately whether \(UR \approx \text{Incidence} \times \text{Duration}\) holds for each group.
    \item Which group has higher unemployment incidence?
    \item Which group has longer unemployment duration?
    \item Explain why youth unemployment rates can be high even when average duration is short.
\end{enumerate}

\item A labour force survey reports:
\begin{itemize}
    \item Working-age population = 5{,}000{,}000
    \item Labour force participation rate = 66\%
    \item Employment rate = 61\%
\end{itemize}

\begin{enumerate}[label=(\alph*)]
    \item Calculate the labour force.
    \item Calculate the number employed.
    \item Calculate the number unemployed.
    \item Calculate the unemployment rate.
\end{enumerate}
\newpage
\item In a recession, the following occurs:
\begin{itemize}
    \item 40{,}000 workers become discouraged and leave the labour force
    \item 30{,}000 workers lose jobs and start actively searching
    \item Initially:
    \begin{itemize}
        \item Labour force = 3{,}000{,}000
        \item Unemployed = 180{,}000
    \end{itemize}
\end{itemize}

\begin{enumerate}[label=(\alph*)]
    \item What is the initial unemployment rate?
    \item What is the new number of unemployed workers?
    \item What is the new labour force?
    \item What is the new unemployment rate?
    \item Briefly explain why the unemployment rate may not fully capture hardship when discouraged workers leave the labour force.
\end{enumerate}

\item Consider the following long-term unemployment data:

\begin{center}
\begin{tabular}{lcc}
\toprule
Country & Total Unemployment Rate & Share Unemployed 12+ Months \\
\midrule
Country X & 6.0\% & 10\% \\
Country Y & 6.0\% & 35\% \\
\bottomrule
\end{tabular}
\end{center}

\begin{enumerate}[label=(\alph*)]
    \item Calculate the long-term unemployment rate for Country X.
    \item Calculate the long-term unemployment rate for Country Y.
    \item Which country has the more serious long-term unemployment problem?
    \item Explain why two countries with the same unemployment rate may have very different labour market conditions.
\end{enumerate}

\item Suppose Canada has:
\begin{itemize}
    \item Unemployment rate = 5.8\%
    \item Labour force participation rate = 65.5\%
    \item Working-age population = 31{,}000{,}000
\end{itemize}

\begin{enumerate}[label=(\alph*)]
    \item Calculate the labour force.
    \item Calculate the number of unemployed workers.
    \item Calculate the number of employed workers.
    \item Calculate the employment rate.
\end{enumerate}

\end{enumerate}

\end{document}
